\scanpage{80}

\deleted{
deux bandes $J\subset J'$, d'où
\[
X_J^*\subset X_{J'}^*
\]
et à regarder le \inserted{pro}voisinage \deleted{pro}tubulaire de $X_J^*$ dans $X_{J'}^*$, on trouve

\margin{On pourrait les noter $V_J^{J'}$.}
\diagram{\begin{tikzcd}[column sep=small,row sep=small,execute at end picture={
\draw ([xshift=-5pt]s80oval.west)
.. controls ([xshift=-5pt]s80oval.north west) and ([yshift=4pt]s80oval.north west) .. ([yshift=4pt]s80oval.north)
.. controls ([yshift=4pt]s80oval.north east) and ([xshift=5pt]s80oval.north east) .. ([xshift=5pt]s80oval.east)
.. controls ([xshift=5pt]s80oval.south east) and ([yshift=-4pt]s80oval.south east) .. ([yshift=-4pt]s80oval.south)
.. controls ([yshift=-4pt]s80oval.south west) and ([xshift=-5pt]s80oval.south west) .. cycle;
}]
& V_{X_J^*,X_{J'}^*}^{\deleted{\unsure{*}}}=\displaystyle\varinjlim_{i,j\in J\atop\deleted{\unsure{i,j\in J'}}}(V_i^{J'},V_{ij}^{J'}) \quad (62)\\
|[alias=s80oval]|\left(X_J^*=\displaystyle\varinjlim_{i,j\in J}(V_i^J,V_{ij}^J)\right)\arrow[ur,"{\substack{(w)\\\sim}}",bend left=12] &
\end{tikzcd}}

(\uline{NB} \unsure{OPS} $J'=I$, \inserted{de plus} $J$ un idéal dans $J'=I$, quitte à \unsure{rapetisser} $I$ en enlevant $J_0^*=\widetilde J\setminus J$; \deleted{$\unsure{\}}$} \deleted{\unsure{Rem.}} \inserted{ou encore,} on est ramené au cas où on a
\[
\unsure{J_0|I_0}\subset J_1\subset J_2
\]
trois idéaux de $I$, où
\[
J=J_1\setminus J_0,\quad J'=J_2\setminus J_0\quad\ldots
\]
).

\margin{On a $V_J^{J'}=V_J^{J_0'}$, si $J\subset J_0'\subset J'$ est formé des $j\in J'$ qui \unsure{majorent|minorent} un élément de $J$. La condition $\unsure{J'|J}=J_0'$ \unsure{signifie} que $X_J^*$ est fermé dans $X_{J'}^*$, car $J$ est un idéal de $J'$.}

Soit Maintenant, de plus, $K\subset I$ une partie de $I$, et considérons
\[
\deleted{V_J^{J'}\unsure{\setminus X_K\cap V_J^{J'}|_{K}}=\unsure{V(X_J^*)}}
\]
\inserted{$(V_J^{J'}\setminus X_{\unsure{K|k}}\cap V_J^{J'}=)$}
\[
V_{J;K}^{J'}=V(X_J^*,X_{J'}^*)\setminus X_K\cap V(X_J^*,X_{J'}^*)
\]
Notons que pour $k\in I$
\[
(63)\qquad V_{J\deleted{\unsure{;}} ,k}^{J'}\overset{\mathrm{déf}}{=}V_J^{J'}\cap X_k=\varnothing
\quad\overset{\inserted{\text{$X_k$ fermé}}}{\Longleftrightarrow}\quad X_J^*\cap X_k=\varnothing
\]
\[
\deleted{\Longleftrightarrow\unsure{X_k\subset\widetilde J|X_k=\varnothing}}
\]
\[
\deleted{\Longleftarrow\forall i\in I,\quad i\leq k\Rightarrow i\notin J}
\]
\deleted{et \unsure{ce}} et comme $X_k$ est la réunion (disjointe) des $X_i^*$ pour $i\leq k$, La relation (63) est conséquence de celle que \deleted{\unsure{si}} $i\leq k\Rightarrow i\notin J$ i.e. du \deleted{fait pour $k$}
}
\mspage{79}
